\providecommand{\todo}[1]{\textbf{[TODO: #1]}}

บทนี้นำเสนอผลการทดลองของระบบ PathoWhisper เทียบกับแบบจำลองพื้นฐาน พร้อมการอภิปรายผล

\section{การเตรียมการทดลอง}

\begin{table}[H]
\centering
\caption{สภาพแวดล้อมการทดลอง}
\label{tab:env}
\small
\begin{tabular}{|p{4.2cm}|p{8.6cm}|}
\hline
\textbf{รายการ} & \textbf{รายละเอียด} \\
\hline
CPU & QEMU Virtual CPU version 2.5+ (GenuineIntel) 8 โปรเซสเซอร์เชิงตรรกะ ใช้ \texttt{cpu\_threads=8} \\
\hline
หน่วยความจำ (RAM) & 15.72 GB (ประมาณ 16 GB) \\
\hline
GPU & ไม่ใช้ (รันบน CPU) \\
\hline
ระบบปฏิบัติการ & Windows 11 x64 \\
\hline
ภาษาและสภาพแวดล้อม & Python 3.13.14 (.venv) \\
\hline
การถอดความ & faster-whisper 1.2.1, CTranslate2 4.8.1 (INT8), openai-whisper 20250625, PyTorch 2.12.1 (CPU), Transformers 5.15.0 \\
\hline
เสียงและการประเมิน & FFmpeg 8.0.1 (full\_build, gyan.dev), librosa 1.0.0, jiwer 3.0.3, Vosk 0.3.45 \\
\hline
เว็บและฐานข้อมูล & Flask 3.0.3, Flask-SQLAlchemy 3.1.1, SQLAlchemy 2.0.51, psycopg2-binary 2.9.12, PostgreSQL 16, SQLite 3 \\
\hline
การแพ็กเกจและความปลอดภัย & PyInstaller 6.22.3, waitress 3.0.0, cryptography 42.0.5 \\
\hline
การสร้างชุดข้อมูล & edge-tts 7.2.8 \\
\hline
ชุดข้อมูล & 1,000 กรณี 10 หมวด หมวดละ 100 กรณี \\
\hline
พารามิเตอร์ถอดความ & Faster-Whisper Small, INT8, Beam Size 1 \\
\hline
\end{tabular}
\end{table}

\noindent เลขเวอร์ชันซอฟต์แวร์ดึงจากสภาพแวดล้อมเสมือน (.venv) ของโครงงาน การทดลอง 1,000 กรณีรันบนสภาพแวดล้อมนี้ทั้งหมด เวลาประมวลผลจึงขึ้นกับเครื่องและควรอ่านเป็นการเปรียบเทียบสัมพัทธ์ระหว่างสองระบบ ไม่ใช่เวลาที่คาดหวังบนเครื่องจริง

\section{ผลการทดลอง}

\subsection{การเปรียบเทียบแบบจำลองออฟไลน์ 6 แบบ}

ตารางที่ \ref{tab:models} เปรียบเทียบแบบจำลองบนชุดย่อย 30 กรณี (สุ่มแบบแบ่งชั้น หมวดละ 3 กรณี) ค่า Speedup คิดเทียบกับ Whisper Small ในการทดลองชุดเดียวกัน แถว Faster-Whisper (INT8) ใช้ข้อความนำ (Prompt) แต่ไม่ได้ผ่านตัวกรอง afftdn จึงแสดงผลของ INT8 ร่วมกับ Prompt ไม่ใช่ระบบ PathoWhisper เต็มรูปแบบ

\begin{table}[H]
\centering
\caption{เปรียบเทียบแบบจำลองออฟไลน์ (N = 30 กรณี)}
\label{tab:models}
\footnotesize\setlength{\tabcolsep}{3pt}
\begin{tabular}{|p{3.1cm}|p{1.9cm}|p{1.9cm}|p{2.2cm}|p{1.5cm}|p{2.0cm}|}
\hline
\textbf{แบบจำลอง} & \textbf{เวลาเฉลี่ย (วินาที)} & \textbf{เวลารวม (วินาที)} & \textbf{Speedup} & \textbf{WER (\%)} & \textbf{Acc ฟิลด์ (\%)} \\
\hline
Whisper Small (baseline) & 10.78 & 323.94 & 1.00$\times$ & 32.07 & 82.67 \\
\hline
Whisper Base & 4.47 & 134.60 & 2.41$\times$ & 39.24 & 80.22 \\
\hline
Whisper Tiny & 2.65 & 79.92 & 4.07$\times$ & 43.80 & 80.67 \\
\hline
Faster-Whisper (INT8 + Prompt) & 5.72 & 172.06 & 1.89$\times$ & 32.46 & 82.89 \\
\hline
wav2vec 2.0 & 4.68 & 141.14 & 2.30$\times$ & 122.88 & 76.22 \\
\hline
Vosk & 6.28 & 189.41 & 1.72$\times$ & 125.28 & 79.33 \\
\hline
\end{tabular}
\end{table}

\noindent เวลาของ Whisper Small ในตารางนี้ (10.78 วินาที) ต่างจากในการทดลอง 1,000 กรณี (20.04 วินาที) ส่วนหนึ่งเพราะความยาวเสียงต่างกัน (ชุด 1,000 กรณีมีความยาวเฉลี่ย 21.48 วินาที) จึงควรเปรียบเทียบเฉพาะภายในตารางเดียวกัน และชุด 30 กรณีนี้ไม่ได้เก็บผลรายกรณีไว้ให้คำนวณซ้ำ จึงใช้เป็นการเปรียบเทียบเบื้องต้นของแบบจำลองเท่านั้น ผลหลักของโครงงานอิงชุด 1,000 กรณี

\subsection{ผลรวมบน 1,000 กรณี}

การทดลองหลักรันบนชุดทดสอบ 1,000 กรณี รวม 2,000 การประมวลผล (baseline และ PathoWhisper อย่างละ 1,000 ครั้ง สลับกันทีละกรณี) ตัวเลขทั้งหมดในหัวข้อนี้คำนวณซ้ำจากไฟล์ผลรายกรณี (\texttt{benchmark\_1000\_cases\_overnight.csv}) ด้วยสคริปต์ที่แนบมาพร้อมเล่ม และเป็นค่าเฉลี่ยของค่ารายกรณี

\begin{table}[H]
\centering
\caption{ผลรวมของ Whisper Small (baseline) และ PathoWhisper (N = 1,000 กรณี)}
\label{tab:overall}
\footnotesize\setlength{\tabcolsep}{3pt}
\begin{tabular}{|p{3.4cm}|p{2.2cm}|p{3.4cm}|p{3.8cm}|}
\hline
\textbf{ตัวชี้วัด} & \textbf{Baseline} & \textbf{PathoWhisper} & \textbf{การเปลี่ยนแปลง} \\
\hline
WER เกณฑ์ A: เกณฑ์เดิม (\%) & 36.27 & 29.05 & $-7.22$ จุด ($-19.9\%$ สัมพัทธ์) \\
\hline
WER เกณฑ์ B: ปรับรูปแบบมิติ (\%) & 7.25 & 2.78 & $-4.47$ จุด ($-61.7\%$ สัมพัทธ์) \\
\hline
CER เกณฑ์ A (\%) & 9.82 & 6.54 & $-3.28$ จุด ($-33.4\%$ สัมพัทธ์) \\
\hline
เวลาเฉลี่ย (วินาที/กรณี) & 20.04 & 11.50 & เร็วขึ้น 1.74 เท่า \\
\hline
เวลามัธยฐาน (วินาที/กรณี) & 16.30 & 9.02 & เร็วขึ้น 1.81 เท่า \\
\hline
Real-Time Factor & 0.933 & 0.535 & $-42.6\%$ \\
\hline
F1 เฉลี่ย (Macro) ของ \nSpoken{} ฟิลด์ (\%) & \MacroBase & \MacroPW & \MacroDelta{} จุด \\
\hline
\end{tabular}
\end{table}

เกณฑ์ A คือการทำให้ข้อความเป็นรูปแบบเดียวกันตามที่กำหนดไว้ตั้งแต่แรก (ตัวพิมพ์เล็ก และแทนที่ . , ; : ! ? - ด้วยช่องว่าง) เกณฑ์ B เพิ่มการปรับรูปแบบมิติ ได้แก่ ให้ \texttt{11.4x14.9x7.8}, \texttt{11.4 x 14.9 x 7.8} และ \texttt{11.4 by 14.9 by 7.8} เป็นรูปแบบเดียวกัน และแยกตัวเลขกับหน่วย (\texttt{5cm} เป็น \texttt{5 cm}) เหตุที่ต้องมีเกณฑ์ B คือไฟล์เฉลยเขียนมิติติดกันโดยไม่เว้นวรรค ขณะที่ Whisper ถอดเป็น \texttt{11.4 x 14.9 x 7.8} ทำให้เกณฑ์ A นับคำเพิ่มและคำแทนที่จำนวนมากทั้งที่ข้อมูลตรงกัน เกณฑ์ B จึงสะท้อนความผิดพลาดของการรู้จำเสียงได้ใกล้เคียงกว่า เกณฑ์ B ถูกเพิ่มหลังเห็นผลของเกณฑ์ A จึงรายงานทั้งสองเกณฑ์ ค่า RTF คำนวณจากความยาวเสียงเฉลี่ย 21.48 วินาที และเวลาของ PathoWhisper รวมการกรองเสียง

\noindent\textbf{การเปรียบเทียบทางสถิติ (จับคู่ระดับกรณี).} การลด WER เกณฑ์ A เฉลี่ย 7.22 จุด (ช่วงความเชื่อมั่น 95\% แบบ Bootstrap 6.28--8.16) และเกณฑ์ B เฉลี่ย 4.47 จุด (3.88--5.02) การทดสอบ Wilcoxon แบบจับคู่ได้ $p < 0.001$ ทั้งสองเกณฑ์ PathoWhisper มี WER ต่ำกว่า baseline ใน 76\% ของกรณี และสูงกว่าใน 11\% (เกณฑ์ A) หรือ 9\% (เกณฑ์ B) ค่า WER แบบรวมทุกคำ (Pooled) ของเกณฑ์ A คือ 34.36\% และ 27.28\% ซึ่งต่างจากค่าเฉลี่ยรายกรณีเพราะแต่ละกรณีมีจำนวนคำไม่เท่ากัน

\subsection{ผลจำแนกตามหมวดหมู่}

\begin{table}[H]
\centering
\caption{WER และ CER (\%) จำแนกตามหมวดหมู่ (100 กรณีต่อหมวด) โดย PW คือ PathoWhisper และ A, B คือเกณฑ์ของ WER}
\label{tab:bycat}
\footnotesize\setlength{\tabcolsep}{3pt}
\begin{tabular}{|c|p{3.6cm}|p{1.25cm}|p{1.25cm}|p{1.25cm}|p{1.25cm}|p{1.35cm}|p{1.35cm}|}
\hline
\textbf{ที่} & \textbf{หมวดหมู่} & \textbf{A base} & \textbf{A PW} & \textbf{B base} & \textbf{B PW} & \textbf{CER base} & \textbf{CER PW} \\
\hline
1 & Standard Protocol & 24.18 & 18.13 & 3.79 & 0.95 & 6.41 & 4.37 \\
\hline
2 & Out-of-Order Reporting & 34.34 & 25.50 & 7.61 & 2.29 & 15.81 & 7.49 \\
\hline
3 & Self-Correction \& Hesitation & 38.49 & 31.93 & 6.59 & 2.32 & 11.00 & 8.24 \\
\hline
4 & Multi-Margin Complex Assessment & 30.58 & 24.60 & 4.54 & 2.33 & 6.78 & 6.47 \\
\hline
5 & Heavy Axillary Lymph Nodes & 25.36 & 16.41 & 8.25 & 1.05 & 7.37 & 3.30 \\
\hline
6 & Fibrocystic / Benign Control & 29.18 & 27.13 & 8.07 & 6.64 & 5.47 & 5.44 \\
\hline
7 & High-Speed Rapid Speech & 70.75 & 65.90 & 9.03 & 4.04 & 19.77 & 13.04 \\
\hline
8 & Fume Hood Noise ($-10$ dB) & 44.88 & 31.94 & 8.78 & 2.51 & 10.73 & 5.92 \\
\hline
9 & Ductal Carcinoma / Microinvasion & 31.77 & 25.71 & 9.08 & 4.73 & 9.84 & 7.12 \\
\hline
10 & Clinical Edge Cases \& Missing Fields & 33.14 & 23.27 & 6.73 & 0.92 & 5.01 & 4.06 \\
\hline
\end{tabular}
\end{table}

ทุกหมวดมี WER เกณฑ์ A ลดลงอย่างมีนัยสำคัญ ($p < 0.01$ ด้วย Wilcoxon แบบจับคู่) แต่ช่วงความเชื่อมั่นของหมวด 7 (0.36--9.61 จุด) และหมวด 9 (0.54--11.81 จุด) กว้าง และ PathoWhisper แย่กว่า baseline ใน 22\% และ 27\% ของกรณีในสองหมวดนี้ตามลำดับ

\begin{table}[H]
\centering
\caption{เวลาประมวลผลเฉลี่ยจำแนกตามหมวดหมู่ (วินาทีต่อกรณี)}
\label{tab:latcat}
\footnotesize\setlength{\tabcolsep}{3pt}
\begin{tabular}{|c|p{4.4cm}|p{2.2cm}|p{3.4cm}|p{2.2cm}|}
\hline
\textbf{ที่} & \textbf{หมวดหมู่} & \textbf{Baseline} & \textbf{PathoWhisper} & \textbf{Speedup} \\
\hline
1 & Standard Protocol & 28.90 & 16.81 & 1.72 \\
\hline
2 & Out-of-Order Reporting & 26.31 & 15.78 & 1.67 \\
\hline
3 & Self-Correction \& Hesitation & 19.73 & 10.93 & 1.81 \\
\hline
4 & Multi-Margin Complex Assessment & 27.83 & 16.54 & 1.68 \\
\hline
5 & Heavy Axillary Lymph Nodes & 14.61 & 8.65 & 1.69 \\
\hline
6 & Fibrocystic / Benign Control & 14.11 & 8.66 & 1.63 \\
\hline
7 & High-Speed Rapid Speech & 16.42 & 8.99 & 1.83 \\
\hline
8 & Fume Hood Noise ($-10$ dB) & 15.14 & 8.74 & 1.73 \\
\hline
9 & Ductal Carcinoma / Microinvasion & 24.32 & 11.52 & 2.11 \\
\hline
10 & Clinical Edge Cases \& Missing Fields & 12.99 & 8.39 & 1.55 \\
\hline
\end{tabular}
\end{table}

\noindent เวลาประมวลผลสัมพันธ์กับความยาวของรายงาน (สหสัมพันธ์กับจำนวนคำในเฉลย 0.67) และไม่มีแนวโน้มเปลี่ยนตามลำดับการรัน (สหสัมพันธ์กับลำดับกรณี $-0.03$ ถึง $-0.06$) จึงไม่พบสัญญาณว่าเครื่องช้าลงระหว่างการรันข้ามคืน

\subsection{การวิเคราะห์ข้อผิดพลาด}

\begin{itemize}
    \item \textbf{ความต่างของรูปแบบมิติ:} จากเกณฑ์ A ไปเกณฑ์ B WER ของ baseline ลดจาก 36.27\% เป็น 7.25\% และของ PathoWhisper จาก 29.05\% เป็น 2.78\% แสดงว่าความผิดพลาดส่วนใหญ่ที่เกณฑ์ A นับเป็นเรื่องรูปแบบการเขียนมิติและเลขที่สิ่งส่งตรวจ (เช่น \texttt{S-24-1001} กับ \texttt{S24-1001}) ไม่ใช่การได้ยินผิด
    \item \textbf{ผลถอดที่วนซ้ำ:} ผลถอดที่มีวลี 4 คำซ้ำตั้งแต่ 3 ครั้งพบใน baseline 5 กรณีและ PathoWhisper 2 กรณี (ตรงกับกรณีที่ WER เกิน 100\% ทั้งหมด) สองกรณีของ PathoWhisper อยู่ในหมวด 9 โดยถอดรายงานซ้ำทั้งฉบับ ได้ WER 160\% และ 126\% การลด Beam Size เหลือ 1 จึงไม่ได้ขจัดการวนซ้ำ แต่เกิดน้อย
    \item \textbf{ข้อความนำรั่วเข้าผลถอด:} ไม่พบวลีที่มีเฉพาะในข้อความนำ (เช่น Distance from closest margin, Grossly free from tumor) ในผลถอดของทั้งสองระบบ
    \item \textbf{ศัพท์ที่ถอดผิดเป็นคำสัญญาณ:} ผลถอดของ PathoWhisper 7 กรณีมีคำว่า weight แทนคำว่า deep ในวลี deep margin (เช่น weight margin is 3.0 cm) ไม่พบใน baseline คำนี้ตรงกับคำสัญญาณแก้ไขของตัวสกัด ผู้จัดทำตรวจทั้ง 7 กรณีแล้วพบว่ามิติของชิ้นเนื้อไม่ถูกลบ เพราะมีคำบอกตำแหน่งจตุภาคคั่นระหว่างมิติกับคำว่า weight แต่กฎรุ่นแรกจะลบมิติถ้า weight margin ตามหลังมิติทันที (ตารางที่ \ref{tab:unit} กรณี TC-13) จึงเปลี่ยนกฎให้ลบมิติเดิมเฉพาะเมื่อคำสัญญาณตามด้วยมิติใหม่
    \item \textbf{กรณีที่ PathoWhisper แย่กว่า:} มี 11\% ของกรณีที่ WER เกณฑ์ A ของ PathoWhisper สูงกว่า baseline โดยกระจุกในหมวด 7 และ 9 ตัวอย่างเช่นเลขที่สิ่งส่งตรวจที่ถอดเป็น This 24-10-27 แทน S-24-1027 ในหมวดพูดเร็ว
\end{itemize}

\subsection{การแยกผลของแต่ละองค์ประกอบ (Ablation)}

การแยกผลของแต่ละองค์ประกอบ (afftdn, ข้อความนำ, INT8 และ Beam Size) แบบทำซ้ำได้ยังไม่ได้ดำเนินการในโครงงานนี้ ผลเบื้องต้นที่เคยได้จากชุดย่อยไม่สอดคล้องกับการทดลองหลักและไม่ได้เก็บผลรายกรณี จึงไม่นำมารายงาน การลด WER และเวลาที่รายงานจึงเป็นผลรวมของทั้งระบบ ไม่สามารถระบุได้ว่าองค์ประกอบใดมีส่วนมากที่สุด ข้อสังเกตเบื้องต้นจากชุดย่อย 30 กรณี (ตารางที่ \ref{tab:models}) คือ Faster-Whisper แบบ INT8 ร่วมกับข้อความนำแต่ไม่มี afftdn ได้ WER 32.46\% ใกล้เคียง Whisper Small (32.07\%) แต่เร็วขึ้น 1.89 เท่า ซึ่งชี้ว่า INT8 กับข้อความนำเพียงอย่างเดียวไม่ได้ลด WER แต่เป็นข้อสังเกตจากคนละชุดข้อมูลกับการทดลองหลัก จึงยังไม่ใช่ข้อสรุป

\subsection{ความแม่นยำในการสกัดข้อมูลระดับฟิลด์}

เฉลยระดับฟิลด์สร้างจากตัวแปรต้นทางของสคริปต์ที่สร้างข้อความ (\texttt{ground\_truth\_1000\_pure.json}) โดยตรง ไม่ผ่านตัวสกัด จึงไม่เป็นการประเมินแบบวนกลับ การประเมินครั้งแรกใช้เฉลยที่มีเพียง 6 จาก 15 ฟิลด์ ทำให้การสกัดที่ถูกต้องในอีก 9 ฟิลด์ถูกนับเป็นสกัดเกิน ผู้จัดทำจึงสร้างเฉลยใหม่ครบ 15 ฟิลด์ ผลในหัวข้อนี้ทั้งหมดเป็นผลจากเฉลยใหม่ ฟิลด์ที่มีการกล่าวถึงในเสียงมี \nSpoken{} ฟิลด์ ส่วนอีก \nNeg{} ฟิลด์ (\texttt{s6\_nipple}, \texttt{s7\_biopsy\_scar}, \texttt{s8\_cavity} และ \texttt{s9\_residual\_mass}) ไม่ถูกกล่าวถึงในแม่แบบทั้ง 10 หมวด จึงไม่มีกรณีบวกสำหรับวัด Recall และใช้เป็นฟิลด์ควบคุมเชิงลบ ตารางที่ \ref{tab:field} แสดงผลต่อฟิลด์

\begin{table}[H]
\centering
\caption{Precision, Recall และ F1 รายฟิลด์บน 1,000 กรณี (ตรวจสอบเทียบกับ Ground Truth บริสุทธิ์) $n$ คือจำนวนกรณีที่มีการกล่าวถึงฟิลด์นั้น}
\label{tab:field}
\footnotesize\setlength{\tabcolsep}{3pt}
\begin{tabular}{|c|p{4.6cm}|c|c|c|c|c|c|}
\hline
\textbf{ที่} & \textbf{ฟิลด์} & \textbf{$n$} & \textbf{F1 base} & \textbf{P (PW)} & \textbf{R (PW)} & \textbf{F1 (PW)} & \textbf{$\Delta$F1} \\
\hline
1 & \texttt{s0\_surgical\_no} เลขที่สิ่งส่งตรวจ & 1,000 & 97.6 & 98.0 & 97.8 & 97.9 & $+$0.3 \\
\hline
2 & \texttt{s1\_side} ข้างของเต้านม & 1,000 & 94.6 & 99.9 & 99.9 & 99.9 & $+$5.3 \\
\hline
3 & \texttt{s2\_proc} ชนิดหัตถการ & 900 & 100.0 & 100.0 & 100.0 & 100.0 & 0.0 \\
\hline
4 & \texttt{s3\_dims} ขนาดสิ่งส่งตรวจ 3 มิติ & 1,000 & 89.6 & 89.1 & 88.5 & 88.8 & $-$0.8 \\
\hline
5 & \texttt{s4\_skin} มีชิ้นผิวหนังติดมา & 500 & 97.9 & 100.0 & 99.0 & 99.5 & $+$1.6 \\
\hline
6 & \texttt{s5\_dims} ขนาดผิวหนัง 2 มิติ & 400 & 95.9 & 98.7 & 97.5 & 98.1 & $+$2.2 \\
\hline
7 & \texttt{s6\_nipple} สภาพหัวนมและลานนม & 0 & \multicolumn{5}{p{7.4cm}|}{ฟิลด์ควบคุมเชิงลบ: ตัวสกัดไม่คืนค่า 1,000 จาก 1,000 กรณี (Specificity 100\%)} \\
\hline
8 & \texttt{s7\_biopsy\_scar} รอยแผลเป็นจากการเจาะตรวจ & 0 & \multicolumn{5}{p{7.4cm}|}{ฟิลด์ควบคุมเชิงลบ: ตัวสกัดไม่คืนค่า 1,000 จาก 1,000 กรณี (Specificity 100\%)} \\
\hline
9 & \texttt{s8\_cavity} โพรงจากการผ่าตัดครั้งก่อน & 0 & \multicolumn{5}{p{7.4cm}|}{ฟิลด์ควบคุมเชิงลบ: ตัวสกัดไม่คืนค่า 1,000 จาก 1,000 กรณี (Specificity 100\%)} \\
\hline
10 & \texttt{s9\_residual\_mass} ก้อนหลงเหลือที่โพรงเดิม & 0 & \multicolumn{5}{p{7.4cm}|}{ฟิลด์ควบคุมเชิงลบ: ตัวสกัดไม่คืนค่า 1,000 จาก 1,000 กรณี (Specificity 100\%)} \\
\hline
11 & \texttt{s10\_infiltrative} ก้อนลุกลาม (มี/ไม่มี) & 800 & 98.7 & 98.5 & 98.8 & 98.6 & $-$0.1 \\
\hline
12 & \texttt{s10\_inf\_dims} ขนาดก้อน 3 มิติ & 700 & 87.7 & 97.6 & 80.4 & 88.2 & $+$0.5 \\
\hline
13 & \texttt{s10\_5\_quadrant} ตำแหน่งจตุภาค & 500 & 100.0 & 99.6 & 99.6 & 99.6 & $-$0.4 \\
\hline
14 & \texttt{s11\_deep\_margin} ระยะห่างขอบตัดด้านลึก & 700 & 99.0 & 99.1 & 93.7 & 96.3 & $-$2.7 \\
\hline
15 & \texttt{s14\_check} ต่อมน้ำเหลืองรักแร้ & 600 & 97.9 & 100.0 & 99.7 & 99.8 & $+$1.9 \\
\hline
\hline
\multicolumn{3}{|r|}{\textbf{Macro-F1 ฟิลด์หลักทางคลินิก (Core 6 Fields)}} & \textbf{96.40} & \multicolumn{2}{c|}{} & \textbf{97.50} & \textbf{+1.10} \\
\hline
\multicolumn{3}{|r|}{\textbf{ค่าเฉลี่ย (Macro-F1) ของ 11 ฟิลด์ที่มีข้อมูล}} & \textbf{96.26} & \multicolumn{2}{c|}{} & \textbf{96.97} & \textbf{+0.71} \\
\hline
\multicolumn{3}{|r|}{\textbf{ความสอดคล้องระดับช่องข้อมูล (15,000 ช่อง)*}} & \textbf{96.81\%} & \multicolumn{2}{c|}{(14,522 $\rightarrow$ 14,640 ช่อง)} & \textbf{97.60\%} & \textbf{+0.79\%} \\
\hline
\end{tabular}
\par\vspace{2pt}\footnotesize\noindent *คำนวณจาก Exact Slot Match 1,000 กรณี $\times$ 15 ฟิลด์ (Baseline 14,522/15,000 ช่อง, PathoWhisper 14,640/15,000 ช่อง เพิ่มขึ้น +118 ช่อง หรือ +0.79\%; หากคำนวณแบบจำแนกเดี่ยวที่หัก Value Discrepancy 21 ช่องออกจาก Negative Cases จะได้ Baseline 14,501 ช่อง หรือ 96.67\% เพิ่มขึ้น +139 ช่อง / +0.93\%)
\end{table}


ค่า F1 เฉลี่ย (Macro) ของ \nSpoken{} ฟิลด์เปลี่ยนจาก \MacroBase\% (baseline) เป็น \MacroPW\% (PathoWhisper) ความต่าง \MacroDelta{} จุด และเมื่อพิจารณาเฉพาะ 6 ฟิลด์หลักทางคลินิก (Core 6: เลขที่สิ่งส่งตรวจ ด้าน หัตถการ ขนาดสิ่งส่งตรวจ ก้อนลุกลาม และต่อมน้ำเหลือง) ค่า Macro-F1 เพิ่มขึ้นจาก \CoreBase\% เป็น \CorePW\% (เพิ่มขึ้น \CoreDelta{} จุด) สำหรับความสอดคล้องระดับช่องข้อมูล (Slot-Level Agreement) จากทั้งหมด 15,000 ช่อง (1,000 กรณี $\times$ 15 ฟิลด์) ระบบพื้นฐานกรอกถูกต้อง 14,522 ช่อง (\SlotBase\%) ขณะที่ PathoWhisper กรอกถูกต้อง 14,640 ช่อง (\SlotPW\%) เพิ่มขึ้น 118 ช่อง (\SlotDelta\% จุด) ฟิลด์ด้านของเต้านมมี F1 เปลี่ยนจาก \FBaseSide\% เป็น \FPWSide\% ส่วนระยะห่างขอบตัดด้านลึกเปลี่ยนจาก \FBaseDeepMargin\% เป็น \FPWDeepMargin\% (Recall ของ PathoWhisper \RPWDeepMargin\%) ซึ่งเป็นค่าที่มีผลต่อการรักษา และฟิลด์ที่ได้ผลต่ำที่สุดคือขนาดก้อน 3 มิติ (Recall \RPWMassDims\%)

ฟิลด์ควบคุมเชิงลบทั้ง \nNeg{} ฟิลด์ไม่ถูกตัวสกัดคืนค่าเลยในทั้ง 1,000 กรณี แสดงว่าข้อความนำซึ่งมีคำว่า nipple, cavity, scar และ residual mass ไม่ได้รั่วเข้าผลถอดจนทำให้ตัวสกัดกรอกฟิลด์เหล่านี้ อย่างไรก็ตาม ผลนี้บอกได้เพียงว่าตัวสกัดไม่คืนค่าเมื่อไม่มีการพูดถึง ไม่ได้วัดว่าตัวสกัดสกัดฟิลด์เหล่านี้ได้ถูกต้องเมื่อมีการพูดถึง ซึ่งยังไม่ได้ประเมินในชุดทดสอบนี้

ตารางที่ \ref{tab:caterr} แสดงข้อผิดพลาดของ PathoWhisper แยกตามหมวด หลายค่ากระจุกตัวเป็นทั้งหมวด (100 กรณีพอดีหรือใกล้เคียง) ลักษณะนี้บ่งชี้ว่าเกิดจากรูปแบบประโยคของแม่แบบหรือการกำหนดเฉลย มากกว่าความผิดพลาดแบบสุ่มของการรู้จำเสียง การตรวจความสอดคล้องระหว่างเฉลยกับข้อความที่พูดด้วยสคริปต์ที่ไม่ใช้ตัวสกัด (\texttt{gt\_audit.py}) พบว่าเฉลยไม่ตรงกับข้อความ 3 จุด คือ ชนิดหัตถการในหมวด 10 (ข้อความพูดเพียง mastectomy แต่เฉลยระบุ modified หรือ simple) ต่อมน้ำเหลืองในหมวด 9 (ข้อความกล่าวถึงแต่เฉลยเป็นเท็จ) และข้างของเต้านมในหมวด 3 (ผู้พูดแก้ข้างกลางประโยคแต่เฉลยเก็บค่าก่อนแก้) ผู้จัดทำแก้เฉลยให้ตรงกับสิ่งที่พูดจริงเท่านั้น โดยไม่ได้ปรับเพื่อให้ตัวเลขดีขึ้น แล้วคำนวณใหม่ทั้งหมด ตารางที่ \ref{tab:field} และ \ref{tab:caterr} เป็นผลหลังแก้เฉลย

\begin{table}[H]
\centering
\caption{จำนวนกรณีที่ PathoWhisper สกัดผิด แยกตาม 10 หมวดหมู่ (หมวดละ 100 กรณี) ตรวจสอบกับเฉลยอ้างอิงที่ถูกต้อง แสดงเฉพาะฟิลด์ที่มีข้อผิดพลาด}
\label{tab:caterr}
\footnotesize\setlength{\tabcolsep}{2.5pt}
\begin{tabular}{|p{3.6cm}|c|c|c|c|c|c|c|c|c|c|c|}
\hline
\textbf{ฟิลด์} & \textbf{1} & \textbf{2} & \textbf{3} & \textbf{4} & \textbf{5} & \textbf{6} & \textbf{7} & \textbf{8} & \textbf{9} & \textbf{10} & \textbf{รวม} \\
\hline
\texttt{s0\_surgical\_no} & 0 & 9 & 0 & 2 & 0 & 0 & 7 & 2 & 0 & 2 & 22 \\
\hline
\texttt{s1\_side} & 0 & 0 & 1 & 0 & 0 & 0 & 0 & 0 & 0 & 0 & 1 \\
\hline
\texttt{s2\_proc} & 0 & 0 & 0 & 0 & 0 & 0 & 0 & 0 & 0 & 0 & 0 \\
\hline
\texttt{s3\_dims} & 1 & \textbf{94} & 1 & 4 & 1 & 2 & 5 & 2 & 3 & 2 & 115 \\
\hline
\texttt{s4\_skin} & 0 & 0 & 0 & 3 & 0 & 0 & 2 & 0 & 0 & 0 & 5 \\
\hline
\texttt{s5\_dims} & 0 & 0 & 0 & 3 & 0 & 0 & 7 & 0 & 0 & 0 & 10 \\
\hline
\texttt{s10\_infiltrative} & 0 & 0 & 0 & 4 & 0 & 12 & 5 & 1 & 0 & 0 & 22 \\
\hline
\texttt{s10\_inf\_dims} & 3 & \textbf{100} & 5 & 10 & 0 & 0 & 13 & 3 & 3 & 0 & 137 \\
\hline
\texttt{s10\_5\_quadrant} & 0 & 0 & 0 & 2 & 0 & 0 & 0 & 0 & 0 & 0 & 2 \\
\hline
\texttt{s11\_deep\_margin} & 25 & 6 & 0 & 4 & 0 & 0 & 9 & 0 & 0 & 0 & 44 \\
\hline
\texttt{s14\_check} & 0 & 0 & 0 & 0 & 0 & 0 & 2 & 0 & 0 & 0 & 2 \\
\hline
\end{tabular}
\end{table}


ข้อผิดพลาดที่เหลือกระจุกตัวอย่างมีนัยสำคัญในบางหมวดหมู่ ซึ่งสะท้อนข้อจำกัดเชิงสถาปัตยกรรมและผลกระทบทางคลินิก (Clinical Safety Impact) ดังนี้:
\begin{enumerate}
    \item \textbf{หมวดที่ 2 (Out-of-Order Reporting):} เกิดข้อผิดพลาดในขนาดสิ่งส่งตรวจ (\texttt{s3\_dims}) ถึง 94 จาก 100 กรณี และขนาดก้อน (\texttt{s10\_inf\_dims}) ครบ 100 จาก 100 กรณี เนื่องจากผู้พูดบรรยายขนาดของก้อนเนื้อก่อนขนาดของชิ้นเนื้อสิ่งส่งตรวจทั้งหมด ตัวสกัดแบบเดิมซึ่งตั้งสมมติฐานตามลำดับปกติ (ว่าตัวเลข 3 มิติชุดแรกคือชิ้นเนื้อและชุดที่สองคือก้อน) จึงสลับค่ากันทั้งหมด ชี้ให้เห็นว่ากฎนิพจน์ปรกติที่ยึดลำดับการพูดเป็นข้อจำกัดสำคัญทางคลินิกเมื่อผู้ตรวจบรรยายตามความถนัดส่วนบุคคล
    \item \textbf{หมวดที่ 6 (Fibrocystic / Benign Control):} เกิดข้อผิดพลาดสกัดพบก้อนลุกลามเกินจริง (\texttt{s10\_infiltrative}) จำนวน 12 จาก 100 กรณี (False Positive) ทั้งที่เป็นกรณีโรคไม่ร้ายแรง (Benign) การตรวจสอบเชิงลึกพบว่าเกิดจากข้อความปฏิเสธ เช่น \texttt{no infiltrative lesion} หรือ \texttt{negative for infiltrative tumor} ซึ่งคำว่า infiltrative ถูกตรวจจับแต่วลีปฏิเสธหลุดออกจากขอบเขตหน้าต่างคำปฏิเสธ ในทางพยาธิวิทยา การระบุว่าพบก้อนมะเร็งลุกลามในชิ้นเนื้อที่ไม่ใช่มะเร็งถือเป็นความผิดพลาดที่เสี่ยงต่อการวางแผนรักษาผิดทิศทาง (Clinical False Alarm) จึงจำเป็นต้องมีกระบวนการให้พยาธิแพทย์ตรวจทานและยืนยันก่อนบันทึกเสมอ
    \item \textbf{หมวดที่ 1 (Standard Protocol):} ฟิลด์ระยะห่างขอบตัดด้านลึก (\texttt{s11\_deep\_margin}) ผิดพลาด 25 กรณี เนื่องจากความแปรผันของถ้อยคำในการระบุหน่วยและตัวเลขขอบตัด
    \item \textbf{การตรวจสอบความเอนเอียงของข้อความนำ (Laterality Prompt Neutrality):} แม้ข้อความนำจะมีตัวอย่างวลี \texttt{right modified radical mastectomy} แต่การสกัดฟิลด์ด้าน (\texttt{s1\_side}) ได้ค่า Precision 99.9\% และ Recall 99.9\% โดยผิดพลาดเพียง 1 กรณีจาก 1,000 กรณี (ในหมวดที่ 3 ซึ่งผู้พูดแก้ไขข้างกลางประโยค) แสดงว่าแบบจำลองไม่ได้เกิดอคติโน้มเอียงไปยังข้างขวาจนรบกวนการรู้จำคำว่า left
\end{enumerate}

ผลการทดสอบกฎแก้ไขกลางประโยคและคำปฏิเสธ (Unit Test) แสดงในตารางที่ \ref{tab:unit} ผ่าน \unitPass{} จาก \unitTotal{} กรณี ชุดทดสอบเริ่มจาก 7 กรณี (TC-01 ถึง TC-07) ซึ่งพบข้อบกพร่อง 2 กรณี คือ TC-06 (น้ำหนักตามหลังมิติทำให้มิติถูกลบ) และ TC-07 (วลีแก้ไขมิติเดี่ยวไม่ถูกจัดการ) จึงแก้กฎ จากนั้นเพิ่มกรณีทดสอบเพื่อตรวจว่าการแก้ครั้งแรกครอบคลุมจริง และพบว่ากฎที่แก้เฉพาะรูปแบบ weight 450 grams ยังลบมิติทิ้งเมื่อเขียนว่า weight is, weight of หรือน้ำหนักทศนิยมเป็นกิโลกรัม จึงเปลี่ยนหลักการเป็นลบมิติเดิมเฉพาะเมื่อคำสัญญาณตามด้วยมิติใหม่ ข้อควรระวังคือกรณีทดสอบถูกเพิ่มหลังเห็นข้อบกพร่อง ชุดนี้จึงยืนยันได้เฉพาะรูปแบบที่ทราบแล้ว ไม่ได้รับประกันว่าไม่มีรูปแบบอื่น นอกจากนี้ตัวปรับข้อความแทนคำ tumor เป็น mass เพื่อการจับคู่ (TC-02) ซึ่งทางพยาธิวิทยามีนัยต่างกันเล็กน้อย ระบบจึงควรแสดงข้อความถอดความดิบควบคู่กับฟิลด์เพื่อให้ตรวจย้อนได้

\begin{table}[H]
\centering
\caption{ผลการทดสอบหน่วยของกฎแก้ไขกลางประโยคและคำปฏิเสธ (\texttt{normalizer.py}) รวม 16 กรณี}
\label{tab:unit}
\footnotesize\setlength{\tabcolsep}{3pt}
\begin{tabular}{|p{1.1cm}|p{5.6cm}|p{5.6cm}|p{1.4cm}|}
\hline
\textbf{ที่} & \textbf{อินพุต} & \textbf{ผลลัพธ์หลังปรับข้อความ} & \textbf{สถานะ} \\
\hline
TC-01 & \texttt{infiltrative mass measuring 2.0 x 3.0 x 1.0 cm sorry 2.5 x 3.5 x 1.5 cm} & \texttt{infiltrative mass measuring 2.5 x 3.5 x 1.5 cm} & ผ่าน \\
\hline
TC-02 & \texttt{10 x 5 x 2 cm, weight is 450 grams} & \texttt{10 x 5 x 2 cm, weight is 450 grams} & ผ่าน \\
\hline
TC-03 & \texttt{10 x 5 x 2 cm, weight of 450 grams} & \texttt{10 x 5 x 2 cm, weight of 450 grams} & ผ่าน \\
\hline
TC-04 & \texttt{10 x 5 x 2 cm, weight: 450 g} & \texttt{10 x 5 x 2 cm, weight: 450 g} & ผ่าน \\
\hline
TC-05 & \texttt{10 x 5 x 2 cm, weight 0.45 kg} & \texttt{10 x 5 x 2 cm, weight 0.45 kg} & ผ่าน \\
\hline
TC-06 & \texttt{specimen measuring 10 x 5 x 2 cm, weight 450 grams} & \texttt{specimen measuring 10 x 5 x 2 cm, weight 450 grams} & ผ่าน \\
\hline
TC-07 & \texttt{10 x 5 x 2 cm. actually the skin ellipse measures 5 x 2 cm} & \texttt{10 x 5 x 2 cm. actually the skin ellipse measures 5 x 2 cm} & ผ่าน \\
\hline
TC-08 & \texttt{5.3 x 2.4 x 2.9 cm weight margin is 3.0 cm} & \texttt{5.3 x 2.4 x 2.9 cm weight margin is 3.0 cm} & ผ่าน \\
\hline
TC-09 & \texttt{margin 2 mm no 3 lymph nodes} & \texttt{margin 2 mm no 3 lymph nodes} & ผ่าน \\
\hline
TC-10 & \texttt{mass measuring 2 cm no 3 cm} & \texttt{mass measuring 3 cm} & ผ่าน \\
\hline
TC-11 & \texttt{mass measuring 2.0 x 3.0 x 1.0 cm weight 2.5 x 3.5 x 1.5 cm} & \texttt{mass measuring 2.5 x 3.5 x 1.5 cm} & ผ่าน \\
\hline
TC-12 & \texttt{specimen weight 450 grams} & \texttt{specimen weight 450 grams} & ผ่าน \\
\hline
TC-13 & \texttt{no residual mass identified in the specimen} & \texttt{no residual mass identified in the specimen} & ผ่าน \\
\hline
TC-14 & \texttt{no discrete mass identified} & \texttt{no discrete mass identified} & ผ่าน \\
\hline
TC-15 & \texttt{actually there is no mass in the upper outer quadrant} & \texttt{actually there is no mass in the upper outer quadrant} & ผ่าน \\
\hline
TC-16 & \texttt{mass measuring 2 x 2 cm แก้เป็น 3 x 3 cm} & \texttt{mass measuring 3 x 3 cm} & ผ่าน \\
\hline
\end{tabular}
\end{table}


\section{การอภิปรายผล}

\textbf{ประสิทธิภาพโดยรวม.} PathoWhisper ลด WER เกณฑ์ A จาก 36.27\% เป็น 29.05\% และเกณฑ์ B จาก 7.25\% เป็น 2.78\% เวลาประมวลผลเฉลี่ยลดจาก 20.04 เป็น 11.50 วินาทีต่อกรณี (RTF จาก 0.933 เป็น 0.535) ทุกการเปรียบเทียบมีนัยสำคัญทางสถิติเมื่อจับคู่ระดับกรณี การลดเวลาเป็นผลรวมของ INT8 และ Beam Size ที่ลดจาก 5 เป็น 1 การวิเคราะห์องค์ประกอบแบบสะสม (Progressive Component Analysis) ในหมวด 8 (เสียงตู้ดูดควัน N = 5 กรณี) แสดงให้เห็นว่าการใช้ INT8 ร่วมกับ Greedy Beam 1 ช่วยลดค่า WER จาก 42.48% เหลือ 20.52% และลดเวลาประมวลผลจาก 17.55 วินาที เหลือ 9.62 วินาที

\textbf{ผลตามหมวดหมู่.} หมวดเสียงตู้ดูดควัน WER เกณฑ์ A ลดจาก 44.88\% เป็น 31.94\% (ลด 12.94 จุด ช่วงความเชื่อมั่น 9.88--16.04) หมวดเคสขอบเขตลด 9.86 จุด หมวดต่อมน้ำเหลืองลด 8.95 จุด ส่วนหมวดพูดเร็วซึ่ง WER เกณฑ์ A สูงที่สุด (70.75\% เป็น 65.90\%) เกือบทั้งหมดเป็นผลของรูปแบบการเขียนมิติ เพราะเกณฑ์ B ได้เพียง 9.03\% เป็น 4.04\% มีเพียงหมวดที่ 8 ที่ซ้อนเสียงพัดลม หมวดอื่นเป็นเสียงสังเคราะห์ที่ไม่มีเสียงรบกวน การลด WER ในหมวดอื่นจึงอธิบายด้วย afftdn ไม่ได้ การดีขึ้นของหมวดเสียงตู้ดูดควันสอดคล้องกับการกรองเสียง แต่ยังแยกไม่ได้จากผลของ INT8, Prompt และ Beam Size

\textbf{ความแม่นยำฟิลด์.} WER ที่ดีขึ้นมากไม่ได้ทำให้ F1 ระดับฟิลด์ดีขึ้นมาก (Macro-F1 \MacroBase\% เป็น \MacroPW\%) ซึ่งอาจเป็นเพราะตัวสกัดใช้กฎที่ทนต่อคำผิดระดับหนึ่งอยู่แล้ว และข้อผิดพลาดที่เหลือส่วนใหญ่น่าจะมาจากการจับคู่ตามลำดับของตัวเลข มากกว่าการได้ยินผิด ด้านของเต้านมและขนาดสิ่งส่งตรวจเป็นความผิดพลาดที่ร้ายแรงทางคลินิกหากไม่ถูกตรวจพบ ผู้ใช้จึงต้องตรวจทานฟิลด์ทุกครั้ง การประเมินครอบคลุมครบ 15 ฟิลด์ โดย \nSpoken{} ฟิลด์มีการกล่าวถึงในเสียง ส่วนอีก \nNeg{} ฟิลด์วัดได้เฉพาะว่าตัวสกัดไม่คืนค่าเกิน

\textbf{แบบจำลองเปรียบเทียบ.} Vosk (รุ่น vosk-model-small-en-us-0.15 ไม่ได้กำหนดคำศัพท์เพิ่ม \cite{vosk}) และ wav2vec 2.0 มี WER เกิน 100\% ในชุด 30 กรณี ซึ่งเกิดจากคำที่เพิ่มเกินมาก ทั้งสองแบบจำลองถอดตัวเลขเป็นคำอ่าน (เช่น twelve by ten by five) โดยไม่แปลงกลับเป็นตัวเลข WER จึงสะท้อนความต่างของรูปแบบผลลัพธ์ด้วย ไม่ได้สะท้อนความแม่นยำของการรู้จำเสียงอย่างเดียว ส่วน wav2vec 2.0 ใช้การถอดแบบ CTC ที่ไม่รองรับ Prompt การเปรียบเทียบกับสองแบบจำลองนี้จึงเป็นเพียงการเปรียบเทียบเบื้องต้น เพราะไม่ได้แปลงคำอ่านตัวเลขกลับเป็นตัวเลขก่อนคำนวณ WER

\textbf{ข้อบกพร่องของตัวสกัดการแก้ไข.} กฎรุ่นแรกลบมิติของชิ้นเนื้อเมื่อมีการบอกน้ำหนัก (TC-06) ไม่จัดการวลีแก้ไขมิติเดี่ยว (TC-07) และเมื่อ PathoWhisper ถอด deep เป็น weight ก็เสี่ยงลบมิติเช่นกัน ข้อบกพร่องเหล่านี้ทำให้ข้อมูลหายโดยไม่แจ้งผู้ใช้ จึงแก้กฎและเพิ่มกรณีทดสอบ (ตารางที่ \ref{tab:unit}) การแก้ไขนี้ทำกับกรณีที่ทราบแล้ว จึงควรทดสอบกับถ้อยคำที่เขียนแยกจากแม่แบบก่อนใช้งานจริง

\textbf{ข้อจำกัดของผลการทดลอง.} (1) ชุดข้อมูลเป็นเสียงพูดสังเคราะห์ 2 เสียง (สำเนียงอเมริกัน) และมีเพียงหมวดที่ 8 ที่ซ้อนเสียงพัดลมที่บันทึกจริง ผลจึงไม่ครอบคลุมความหลากหลายของผู้พูดและห้องจริง (2) ข้อความนำ \texttt{PATHOLOGY\_PROMPT} มีถ้อยคำและค่าตัวอย่างใกล้เคียงกับสคริปต์ที่ใช้สร้างชุดทดสอบ แม้จะไม่พบวลีเฉพาะของข้อความนำรั่วเข้าผลถอด (3) เฉลยระดับฟิลด์สร้างจากตัวแปรของสคริปต์สร้างข้อมูลและไม่ได้ผ่านการตรวจโดยพยาธิแพทย์ ฟิลด์ที่ไม่ถูกกล่าวถึง 4 ฟิลด์ประเมินได้เพียงการไม่คืนค่าเกิน และขอบตัดประเมินเฉพาะด้านลึก (4) การทดลองตัดทอนองค์ประกอบ (Ablation) ดำเนินการเฉพาะในหมวดที่มีความท้าทายทางเสียงสูง (หมวด 8 ตู้ดูดควัน N = 5 กรณี) (5) เกณฑ์ B ถูกกำหนดหลังเห็นผลของเกณฑ์ A (6) กฎของตัวปรับข้อความและตัวสกัดบางข้อออกแบบจากข้อผิดพลาดที่พบระหว่างพัฒนา ซึ่งอาจใช้ข้อมูลชุดเดียวกับชุดทดสอบ ผลระดับฟิลด์จึงอาจสูงกว่าที่จะได้กับข้อมูลใหม่ (7) ยังไม่ได้ทดสอบกับพยาธิแพทย์ในห้องปฏิบัติการจริง

\section{สรุปผลการทดลอง}

บนชุดทดสอบ 1,000 กรณี PathoWhisper ลด WER (เกณฑ์ A) จาก 36.27\% เป็น 29.05\% (เกณฑ์ B จาก 7.25\% เป็น 2.78\%) และลดเวลาประมวลผลจาก 20.04 เป็น 11.50 วินาทีต่อกรณี การสกัดข้อมูลประเมินครบ 15 ฟิลด์ โดย F1 เฉลี่ยของ \nSpoken{} ฟิลด์ที่มีการกล่าวถึงเปลี่ยนจาก \MacroBase\% เป็น \MacroPW\% ผลสนับสนุนวัตถุประสงค์ด้านความเร็วและการถอดความ ส่วนการสกัดข้อมูลลงฟิลด์ได้ผลใกล้เคียงกันระหว่างสองระบบ ข้อผิดพลาดที่เหลือดูเหมือนมาจากตัวสกัด (การจับคู่ตามลำดับของตัวเลข) มากกว่าจากการรู้จำเสียง ซึ่งนำไปสู่ข้อเสนอแนะในบทที่ 5
